% LaTeX report generated by Copilot assistant
\documentclass[12pt,a4paper]{article}
\usepackage{fontspec}
\usepackage{polyglossia}
\setdefaultlanguage{vietnamese}
\setotherlanguage{english}
\newfontfamily\vnfont{Times New Roman}
\setmainfont{Times New Roman}
\usepackage{hyperref}
\usepackage{listings}
\usepackage{xcolor}

% configure listings for various languages
\lstset{
    basicstyle=\ttfamily\small,
    breaklines=true,
    frame=single,
    numbers=left,
    numberstyle=\tiny,
    keywordstyle=\color{blue},
    commentstyle=\color{gray},
    stringstyle=\color{red},
    showstringspaces=false,
}

\title{T\'A?I LI\?U B\'A?O C\'A?O PROJECT\\
       \large{\textit{PROJECT REPORT DOCUMENTATION}}}
\author{Sinh vi\^en: (c?n x\'ac nh?n)}
\date{\today}

\begin{document}

\maketitle
\thispagestyle{empty}
\newpage

\tableofcontents
\newpage

\section{Th\ông tin Project}
\begin{itemize}
    \item T\^en project: Asis03 (FE React + BE Spring Boot)
    \item Sinh vi\^en: (c\'an x\'ac nh?n t\^en, m\'a s\'inh vi\^en)
    \item M\'o t\'a ng\'an: Ứng dụng quản lý đặt phòng khách sạn với phân quyền nhân viên và khách hàng,
    sử dụng React cho phần giao diện và Spring Boot cho API.
    \item Công nghệ chính: React, Vite, Axios, React Router, Context API (FE); Spring Boot,
    Spring Security, JPA/Hibernate, MySQL (BE).
    \item Link GitHub: \url{https://github.com/maidanghuy/SBA301_SP26} (hiện tại branch main)
\end{itemize}

\section*{A. C\'au tr\'uc Project FE}
\addcontentsline{toc}{section}{A. C\'au tr\'uc Project FE}
\subsection{C\'au tr\'uc Project}
C\^u tr\'uc thư mục \\texttt{src/} như sau:
\begin{verbatim}
src/
  App.jsx
  api/
    axiosClient.js
    authService.js
    bookingService.js
    customerService.js
    roomService.js
    roomTypeService.js
  assets/
  components/
    common/
      AdminMenu.jsx
      Footer.jsx
      Header.jsx
  context/
  data/
  hooks/
    useAuth.js
  layouts/
    AdminLayout.jsx
    MainLayout.jsx
  pages/
    NotFound.jsx
    admin/
      BookingsAdmin.jsx
      CustomersAdmin.jsx
      Dashboard.jsx
      RoomsAdmin.jsx
    Auth/
      Login.jsx
      Register.jsx
    Customer/
      MyBookings.jsx
      NewBooking.jsx
      Profile.jsx
    Home/
      Home.jsx
    Rooms/
      Rooms.jsx
  routes/
    ProtectedRoute.jsx
    RoleRoute.jsx
  store/
    auth/
      AuthContext.jsx
      AuthProvider.jsx
  styles/
    index.css
\end{verbatim}
Mỗi thư mục có chức năng:
\begin{itemize}
    \item \texttt{api/}: chứa các service gọi API bằng axios.
    \item \texttt{components/}: các component tái sử dụng (header, footer, menu...).
    \item \texttt{hooks/}: custom hooks, ví dụ \texttt{useAuth} quản lý trạng thái xác thực.
    \item \texttt{layouts/}: bố cục chính cho khách hàng (MainLayout) và quản trị (AdminLayout).
    \item \texttt{pages/}: các trang theo từng luồng nghiệp vụ.
    \item \texttt{routes/}: định nghĩa route bảo vệ và phân quyền.
    \item \texttt{store/}: context/provider cho xác thực.
    \item \texttt{styles/}: CSS chung.
\end{itemize}

\subsection{C\'a?i \d?\'at & C\'au h\'inh Axios}
\subsubsection{C\'a?i \d?\'at axios}
Dự án sử dụng npm/yarn để cài \texttt{axios}. Lệnh tham khảo:
\begin{verbatim}
npm install axios
\end{verbatim}

\subsubsection{C\'au h\'inh \texttt{axiosClient.js}}
File tại \texttt{src/api/axiosClient.js}:
\begin{lstlisting}[language=JavaScript]
import axios from "axios";

const axiosClient = axios.create({
    baseURL: "http://localhost:8081/api",
    headers: { "Content-Type": "application/json" },
});

// attach access token to every request
axiosClient.interceptors.request.use((config) => {
    const token = localStorage.getItem("access_token");
    if (token) config.headers.Authorization = `Bearer ${token}`;
    return config;
});

// helper for subscribers waiting during token refresh
let isRefreshing = false;
let refreshSubscribers = [];
const subscribeTokenRefresh = (cb) => {
    refreshSubscribers.push(cb);
};
const onRefreshed = (newToken) => {
    refreshSubscribers.forEach((cb) => cb(newToken));
    refreshSubscribers = [];
};

// create a plain axios instance for refresh requests (no interceptors)
export const refreshClient = axios.create({
    baseURL: axiosClient.defaults.baseURL,
    headers: { "Content-Type": "application/json" },
});

// simple response interceptor to unwrap payloads like { status,message,data:... }
axiosClient.interceptors.response.use(
    (res) => res.data.data,
    async (err) => { /* omitted for brevity */ }
);

export default axiosClient;
\end{lstlisting}
Phần còn lại xử lý tự động refresh token và thông báo lỗi.

\subsubsection{T\'ao service gọi API}
Danh sách file service trong \texttt{src/api}:
\begin{itemize}
    \item \texttt{authService.js}: đăng nhập, đăng ký, lấy thông tin \texttt{me},
    làm mới token.
    \item \texttt{bookingService.js}: CRUD đặt phòng, lịch sử, hủy, danh sách cho admin.
    \item \texttt{customerService.js}: (n\'ean) các API liên quan khách hàng.
    \item \texttt{roomService.js}: APIs cho khách hàng xem phòng.
    \item \texttt{roomTypeService.js}: loại phòng.
\end{itemize}
Mỗi service xuất một đối tượng chứa các phương thức map đến endpoint.

\subsection{B\'ao v\'e router}
Hai component router bảo vệ:
\subsubsection{ProtectedRoute}
\begin{lstlisting}[language=JavaScript]
import { Navigate, Outlet } from "react-router-dom";
import useAuth from "../hooks/useAuth";

export default function ProtectedRoute() {
  const { isAuthenticated } = useAuth();
  if (!isAuthenticated) return <Navigate to="/login" replace />;
  return <Outlet />;
}
\end{lstlisting}
Chỉ cho phép người dùng đã đăng nhập truy cập.

\subsubsection{RoleRoute}
\begin{lstlisting}[language=JavaScript]
import { Navigate, Outlet } from "react-router-dom";
import useAuth from "../hooks/useAuth";

export default function RoleRoute({ allow = [] }) {
  const { user } = useAuth();
  const role = user?.role; // "STAFF" | "CUSTOMER"
  if (!role) return <Navigate to="/login" replace />;
  if (allow.length > 0 && !allow.includes(role))
    return <Navigate to="/" replace />;
  return <Outlet />;
}
\end{lstlisting}
Giới hạn theo vai trò.

\subsection{C\'ap nhật UI & CRUD}
\subsubsection{Category READ}
(Repository không có mục Category, cần xác nhận.)
\subsubsection{Category CREATE/UPDATE}
(như tr\'en)
\subsubsection{News READ}
(không tìm thấy cấu trúc tin tức – có thể không áp dụng)
\subsubsection{News CREATE/UPDATE}
(không có)
\subsubsection{Users READ}
Trang quản trị khách hàng: \texttt{CustomersAdmin.jsx} liệt kê người dùng.
\subsubsection{Users CREATE/UPDATE}
Trang này chỉ xử lý xóa và cập nhật trạng thái; chi tiết triển khai trong mã.

\subsection{C\'ap nhật \texttt{App.jsx}}
Routing chính như sau:
\begin{lstlisting}[language=JavaScript]
// (đã tr\'ich từ repo)
// ... xem phần code t?i tr\^en
\end{lstlisting}
- \texttt{MainLayout} chứa các tuyến public: \texttt{"/"}, \texttt{"/rooms"},
  đăng nhập/đăng ký.
- Các route bên trong \texttt{ProtectedRoute} yêu cầu xác thực.
- \texttt{RoleRoute allow={["CUSTOMER"]}} cung cấp luồng khách hàng cho
  trang cá nhân, đặt phòng.
- \texttt{RoleRoute allow={["STAFF"]}} chuyển sang \texttt{AdminLayout}
  chứa Dashboard và các CRUD quản lý.

\subsection{Kết quả Login + CRUD}
\begin{itemize}
    \item Mục chứa hình minh họa (nếu có) được đặt trong thư mục
    \texttt{docs/images} hoặc \texttt{docs/screenshots}; hiện tại không tồn tại,
    do đó thêm placeholders:
    \begin{figure}[h]
    \centering
    (Chèn hình ở đây: màn hình đăng nhập)
    \caption{Login}
    \end{figure}
    \begin{figure}[h]
    \centering
    (Chèn hình ở đây: CRUD Category)
    \caption{CRUD Category}
    \end{figure}
    \begin{figure}[h]
    \centering
    (Chèn hình ở đây: CRUD News)
    \caption{CRUD News}
    \end{figure}
    \begin{figure}[h]
    \centering
    (Chèn hình ở đây: CRUD Users)
    \caption{CRUD Users}
    \end{figure}
\end{itemize}

\newpage
\section*{B. C\'au tr\'uc Project BE}
\addcontentsline{toc}{section}{B. C\'au tr\'uc Project BE}
\subsection{C\'au tr\'uc Project}
C\^u tr\'uc l\^a thư mục tiêu chuẩn Maven/Spring Boot:
\begin{verbatim}
src/main/java/org/example/asis03/
  Asis03Application.java
  common/
  config/
  controller/
    AuthController.java
    RoomPublicController.java
    AdminRoomController.java
    AdminRoomTypeController.java
    AdminCustomerController.java
    AdminBookingController.java
    BookingController.java
    UserController.java
  dto/
    request/ ...
    response/ ...
  entity/
    AS03Customer.java
    AS03RoomInformation.java
    AS03RoomType.java
    AS03BookingReservation.java
    AS03BookingDetail.java
    AS03BookingDetailId.java
    BookingStatus.java
    Role.java
    BaseEntity.java
  exception/
  repository/
  service/
    AuthService.java
    RoomInformationService.java
    ...
  service/impl/
    AuthServiceImpl.java
    RoomInformationServiceImpl.java
    ...

src/main/resources/application.yaml

test/... (có một lớp \texttt{Asis03ApplicationTests.java})

target/ ... (thư mục build)

db.json (FE json file)
\end{verbatim}

\subsection{\texttt{Asis03Application.java}}
Tệp khởi động Spring Boot, chứa hàm \texttt{main}. Thiết lập context,
quét các \texttt{@Component} và khởi tạo server.

\subsection{Các controller}
Danh sách controller và endpoint chính:
\begin{itemize}
    \item \texttt{AuthController} (\texttt{/api/auth}): register, login,
    me, refresh token.
    \item \texttt{RoomPublicController} (\texttt{/api/rooms}): xem danh sách
    phòng và lấy theo id cho người dùng chung.
    \item \texttt{AdminRoomController} (\texttt{/api/admin/rooms}): CRUD phòng
    dành cho STAFF (có \texttt{@PreAuthorize}).
    \item \texttt{AdminRoomTypeController} (\texttt{/api/admin/room-types}): quản lý
    loại phòng.
    \item \texttt{AdminCustomerController} (\texttt{/api/admin/customers}): xem,
    cập nhật khách hàng.
    \item \texttt{AdminBookingController} (\texttt{/api/admin/bookings}): danh sách
    và thao tác đặt phòng.
    \item \texttt{BookingController} (\texttt{/api/bookings}): nghiệp vụ đặt,
    hủy, lịch sử.
    \item \texttt{UserController} (\texttt{/api/users}): (nếu có) các endpoint
    chung cho người dùng.
\end{itemize}

\subsection{Service + Service Impl}
Các interface trong \texttt{service/} định nghĩa nghiệp vụ. Impl xử lý logic:
như xác thực, truy vấn repository, thao tác với entity.
Luồng điển hình: controller gọi \texttt{roomService.getAllAdmin()} ->
impl với \texttt{roomRepository.findAll()} -> map sang DTO -> trả lại.
Các dịch vụ còn có \texttt{AuthServiceImpl} để xử lý đăng nhập,
refresh token, đăng ký. Logging, kiểm tra hợp lệ xảy ra tại đây.
Các dịch vụ admin tách biệt \texttt{Admin*Service}.

\subsection{Entity}
Danh sách entity chính:
\begin{itemize}
    \item \texttt{AS03Customer}: thông tin khách hàng, quan hệ với
    đặt phòng (1-n).
    \item \texttt{AS03RoomInformation}: phòng, thuộc loại phòng, có trạng thái.
    \item \texttt{AS03RoomType}: loại phòng (ví dụ Single, Double).
    \item \texttt{AS03BookingReservation}: bảng đặt phòng chính, liên kết tới
    khách hàng.
    \item \texttt{AS03BookingDetail}, \texttt{AS03BookingDetailId}: chi tiết
    mỗi phòng trong một đặt phòng (composite key).
    \item \texttt{BookingStatus}: enum trạng thái đặt (PENDING, CONFIRMED, ...)
    \item \texttt{Role}: enum vai trò người dùng (CUSTOMER, STAFF).
    \item \texttt{BaseEntity}: lớp cha chứa ID, timestamp tạo/sửa.
\end{itemize}
Quan hệ chính
\texttt{AS03BookingReservation} 1-n \texttt{AS03BookingDetail},
\texttt{AS03RoomInformation} n-1 \texttt{AS03RoomType},
\texttt{AS03Customer} 1-n \texttt{AS03BookingReservation}.

\end{document}
